% STATUS: ACTIVE | NOTE ON THE PREDICT LINE OF THE GAIN: EULER STEP, ITS RATCHET,
%   THE EXACT STEP, AND WHERE THE ESTIMATE ENTERS (2026-09-02)
%
%   One line of the filter is the subject:  a_hat[k+1] = a_hat[k] + slope x step.
%   S1  the constant-b law and its curvature (a'' = -2 b^2 (1-a)^3)
%   S2  one step of the law: truth vs the Euler step; the difference is
%       +b^2 (1-a)^3 (step)^2, same sign for either direction of travel (ratchet)
%   S3  the ratchet summed along the two canonical trajectories (open-loop,
%       zero-parameter): canon 1 Hz +0.0066 (measured closed-loop zero-noise
%       +0.0058), Meng cosine descent +0.0001 (measured +0.00003)
%   S4  the exact step: 1/(1-a) is affine in w with slope b, so the step in
%       1/(1-a) has no truncation at any order in the step size
%   S5  where the estimate a_hat enters the predict line, per arm and per step
%       scheme, and the second-order mean rule E[f(x)] = f(E x) + 1/2 f'' Var x.
%       The exact step as implemented for the exogenous-slope oracle arm
%       (b_eff = a'_true(w_d)/(1-a_hat)^2) puts a_hat into a denominator; the
%       Euler oracle arm has no a_hat in its increment at all. That, and only
%       that, is the "cost" of the exact step: it exists for the oracle arm,
%       not for production, where a_hat is in the increment under Euler too.
%   S6  an exact step for an exogenous slope that contains no a_hat.
%   Numbers: b = 8/9, Ts = 0.625 ms; canon deep band [1.10, 3.32], Meng
%   descent 6.667 -> 1.111 in 10 s (cosine) + 2 s hold. Bias numbers quoted as
%   (a_hat - a_bar), zero-noise, 2026-09-01 (run_arms_seedtruth_pair.m).
%   Companion: 0902_formC_aptrue_4state.tex (full 4-state derivation of the
%   oracle arm with the exact step and the second-order mean term S5/S6 there).
% ======================================================
% Compile: xelatex 0902_law_step_euler_vs_exact.tex
% Path:    reference/eq17_analysis/derivation/0902_law_step_euler_vs_exact.tex
% ======================================================

\documentclass[11pt,a4paper,fontset=none]{ctexart}
\usepackage[fleqn]{amsmath}
\usepackage{amssymb,bm}
\usepackage{graphicx}
\setlength{\mathindent}{0pt}
\usepackage{geometry}
\geometry{margin=2.0cm}

\setlength{\parindent}{0pt}
\setlength{\parskip}{0.80em}
\linespread{1.08}
\setlength{\abovedisplayskip}{10pt}
\setlength{\belowdisplayskip}{10pt}

\newcommand{\ba}{\bar{a}}
\newcommand{\bw}{\bar{w}}
\newcommand{\E}{\mathbb{E}}

\begin{document}

\section{Gain law and its curvature}

(bars denote normalized quantities: lengths by $R$, the gain by
$a_{\mathrm{nom}} = \Delta t/\gamma_N$, $\ba = a/a_{\mathrm{nom}}$,
$0 \le \ba \le 1$, far field $1$, contact $0$; $\bw$ is the height in units of $R$;
$\ba' = d\ba/d\bw$, $\ba'' = d^{2}\ba/d\bw^{2}$;
hats sit on estimates; $e_{\theta} = \theta - \hat{\theta}$ as in the SSOT,
but every bias number below is quoted as $\hat{\ba} - \ba$;
$\E[\cdot]$ is the mean over noise realizations)

\begin{equation*}
\ba' = b\,(1-\ba)^{2}
\qquad\text{(the law: the slope is a function of the gain itself, } b \text{ constant)}
\end{equation*}

\begin{align*}
\ba'' &= \frac{d\ba'}{d\bw}
      = \frac{d\ba'}{d\ba}\,\frac{d\ba}{d\bw}
      = \bigl[-2b\,(1-\ba)\bigr]\bigl[b\,(1-\ba)^{2}\bigr]
      = -2\,b^{2}\,(1-\ba)^{3} \\
\ba''' &= \frac{d\ba''}{d\bw}
      = \bigl[-2b^{2}\cdot 3\,(1-\ba)^{2}\cdot(-1)\bigr]\bigl[b\,(1-\ba)^{2}\bigr]
      = 6\,b^{3}\,(1-\ba)^{4}
\end{align*}

($\ba'' < 0$ everywhere: the curve bends downward, steep at the wall and flat
far away. Values on the published $c_\perp$ curve, $b = 8/9$)

\begin{center}
\begin{tabular}{lcccc}
$\bw$ & $\ba$ & $\ba' = b(1-\ba)^{2}$ & $\ba'' = -2b^{2}(1-\ba)^{3}$ & $(1-\ba)^{3}$ \\
\hline
$1.10$ (trough) & $0.087$ & $0.740$ & $-1.20$   & $0.760$ \\
$2.00$          & $0.471$ & $0.249$ & $-0.234$  & $0.148$ \\
$6.67$ (start)  & $0.833$ & $0.025$ & $-0.0074$ & $0.0047$
\end{tabular}
\end{center}

\section{One step of the law: truth versus the Euler step}

(the height moves from $\bw[k]$ to $\bw[k+1]$; the true gain moves along the
curve; the Euler step moves along the tangent at $\bw[k]$. Both start from the
same point, $\hat{\ba}[k] = \ba[k]$, so the difference is the truncation of
this single step)

\begin{align*}
\ba[k+1] &= \ba[k]
  + \ba'\bigl(\bw[k]\bigr)\bigl(\bw[k+1]-\bw[k]\bigr)
  + \tfrac{1}{2}\,\ba''\bigl(\bw[k]\bigr)\bigl(\bw[k+1]-\bw[k]\bigr)^{2} \\
  &\qquad
  + \tfrac{1}{6}\,\ba'''\bigl(\bw[k]\bigr)\bigl(\bw[k+1]-\bw[k]\bigr)^{3}
  + \cdots \\
\hat{\ba}[k+1] &= \hat{\ba}[k]
  + b\,\bigl(1-\hat{\ba}[k]\bigr)^{2}\bigl(\bw[k+1]-\bw[k]\bigr)
\qquad\text{(Euler)}
\end{align*}

\begin{align*}
\hat{\ba}[k+1] - \ba[k+1]
  &= -\tfrac{1}{2}\,\ba''\bigl(\bw[k]\bigr)\bigl(\bw[k+1]-\bw[k]\bigr)^{2}
     -\tfrac{1}{6}\,\ba'''\bigl(\bw[k]\bigr)\bigl(\bw[k+1]-\bw[k]\bigr)^{3}
     - \cdots \\
  &= +\,b^{2}\bigl(1-\ba[k]\bigr)^{3}\bigl(\bw[k+1]-\bw[k]\bigr)^{2}
     -\,b^{3}\bigl(1-\ba[k]\bigr)^{4}\bigl(\bw[k+1]-\bw[k]\bigr)^{3}
     + \cdots
\end{align*}

(the leading term is $+b^{2}(1-\ba)^{3}(\bw[k+1]-\bw[k])^{2}$: positive for a
step toward the wall, $\bw[k+1]-\bw[k] < 0$, and positive for a step away from
it, $\bw[k+1]-\bw[k] > 0$, because it is the square of the step. The tangent
line lies above a downward-bending curve on both sides of the tangent point,
so the Euler estimate lands above the truth in either direction)

\begin{align*}
\hat{\ba} - \ba &= b^{2}(1-\ba)^{3}(\bw[k+1]-\bw[k])^{2}
                 + b^{2}(1-\ba)^{3}(\bw[k+1]-\bw[k])^{2}
                 + \mathcal{O}\bigl((\bw[k+1]-\bw[k])^{3}\bigr) \\
                &= 2\,b^{2}(1-\ba)^{3}(\bw[k+1]-\bw[k])^{2}
\qquad\text{(one step toward the wall, the same step back)}
\end{align*}

(nothing cancels over a round trip; the error only accumulates. The cubic term
$-b^{3}(1-\ba)^{4}(\bw[k+1]-\bw[k])^{3}$ changes sign with the direction and
does cancel over a round trip; the ratchet is the even power)

\section{The ratchet summed along a trajectory (open loop)}

\begin{equation*}
\hat{\ba}[N] - \ba[N]
  = \sum_{k=0}^{N-1} b^{2}\bigl(1-\ba_{\mathrm{true}}(\bw_d[k])\bigr)^{3}
    \bigl(\Delta\bw_d[k]\bigr)^{2}
\qquad\text{(command path, no noise, no filter feedback)}
\end{equation*}

(two factors must be large in the same step: $(1-\ba)^{3}$ is large only near
the wall, $(\Delta\bw_d)^{2}$ only where the command moves fast)

\begin{center}\footnotesize\setlength{\tabcolsep}{4pt}
\begin{tabular}{lccccc}
trajectory & $\max|\Delta\bw_d|$ & steps $|\Delta\bw_d| > 10^{-3}$
  & open-loop sum & share $\bw_d < 2$ & measured, zero noise \\
\hline
canon 1 Hz in $[1.10,\,3.32]$ & $2.07\times10^{-2}$ & 4278 & $+0.0066$ & $0.58$ & $+0.0058$ \\
Meng cosine $6.67 \to 1.11$   & $5.45\times10^{-4}$ & 0    & $+0.0001$ & $0.51$ & $+0.00003$
\end{tabular}
\end{center}

(canon: the oscillation carries the particle through $\bw < 2$ at full speed
five times, peak step $2\pi\cdot 1\,\mathrm{Hz}\cdot 1.11\,R\cdot 0.625\,\mathrm{ms}
= 4.4\times10^{-3}$, where $(1-\ba)^{3} \approx 0.5$. Meng: the cosine profile
decelerates to zero at the trough, so $(\Delta\bw_d)^{2}$ is small exactly where
$(1-\ba)^{3}$ is large, and where the command is fast, $(1-\ba)^{3} \approx 0.005$.
The 12\% gap between the open-loop sum and the measured closed-loop value is the
$y_1$ update pulling part of the injected error back; it does not change the sign)

\section{The exact step}

\begin{equation*}
\frac{d}{d\bw}\Bigl[\frac{1}{1-\ba}\Bigr]
  = \frac{\ba'}{(1-\ba)^{2}}
  = \frac{b\,(1-\ba)^{2}}{(1-\ba)^{2}}
  = b
\qquad\text{(a constant: } 1/(1-\ba) \text{ is a straight line in } \bw\text{)}
\end{equation*}

\begin{equation*}
\frac{1}{1-\ba[k+1]} = \frac{1}{1-\ba[k]} + b\,\bigl(\bw[k+1]-\bw[k]\bigr)
\qquad\text{(exact for any step size, no term is dropped)}
\end{equation*}

\begin{align*}
\ba[k+1] &= 1 - \frac{1}{\dfrac{1}{1-\ba[k]} + b\,\bigl(\bw[k+1]-\bw[k]\bigr)} \\
         &= \ba[k]
  + \frac{b\,(1-\ba[k])^{2}\,\bigl(\bw[k+1]-\bw[k]\bigr)}
         {1 + b\,(1-\ba[k])\,\bigl(\bw[k+1]-\bw[k]\bigr)}
\end{align*}

(expand the exact step in the step size and compare with the Taylor series of
the truth in Section 2, term by term)

\begin{align*}
\Delta &:= \bw[k+1]-\bw[k] \\
\frac{b\,(1-\ba)^{2}\,\Delta}{1 + b\,(1-\ba)\,\Delta}
  &= b\,(1-\ba)^{2}\,\Delta
   - b^{2}(1-\ba)^{3}\,\Delta^{2}
   + b^{3}(1-\ba)^{4}\,\Delta^{3}
   - \cdots \\
\ba'\,\Delta + \tfrac{1}{2}\ba''\,\Delta^{2} + \tfrac{1}{6}\ba'''\,\Delta^{3} + \cdots
  &= b\,(1-\ba)^{2}\,\Delta
   - b^{2}(1-\ba)^{3}\,\Delta^{2}
   + b^{3}(1-\ba)^{4}\,\Delta^{3}
   + \cdots
\end{align*}

(identical at every order: the exact step is the whole series of the
constant-$b$ law, so the ratchet of Section 2 is zero at any step size.
The Euler step is this series cut after the first term. What the exact step
does not remove is the difference between the law and the truth,
$b_{\mathrm{true}}(\ba) \ne b$; that is a model error, not a step-size error,
and it is the same at every step size. Measured: canon zero-noise
$+0.0058 \to -0.00001$, Meng $+0.00003 \to -0.00008$)

\section{Where the estimate enters the predict line}

(the filter's height increment is
$\Delta\bw_d[k] + (1-\lambda_c)\,\delta\hat{\bw}_3[k]$:
the commanded move plus the part of the estimated disturbance the model expects
to relax this step. In each arm the predict line is written out in full; the
question asked of each line is only whether it is linear in the quantities that
fluctuate across noise realizations, $\hat{\ba}$ and $\delta\hat{\bw}_3$)

\textbf{Oracle arm, slope exogenous at the commanded height, Euler}

\begin{equation*}
\hat{\ba}[k+1] = \hat{\ba}[k]
  + \ba'_{\mathrm{true}}\bigl(\bw_d[k]\bigr)
    \bigl[\Delta\bw_d[k] + (1-\lambda_c)\,\delta\hat{\bw}_3[k]\bigr]
\end{equation*}

($\ba'_{\mathrm{true}}(\bw_d[k])$ is a known number at step $k$;
$\hat{\ba}[k]$ enters once, to the first power; $\delta\hat{\bw}_3[k]$ enters
once, to the first power. The line is linear. Its mean is the line evaluated at
the means, whatever the noise. Measured bias $+0.0001$, Meng, seed at truth)

\textbf{Oracle arm, slope exogenous, exact step as implemented
($b_{\mathrm{eff}} = \ba'_{\mathrm{true}}(\bw_d[k])/(1-\hat{\ba}[k])^{2}$)}

\begin{align*}
\frac{1}{1-\hat{\ba}[k+1]}
  &= \frac{1}{1-\hat{\ba}[k]}
   + \frac{\ba'_{\mathrm{true}}\bigl(\bw_d[k]\bigr)}{\bigl(1-\hat{\ba}[k]\bigr)^{2}}
     \bigl[\Delta\bw_d[k] + (1-\lambda_c)\,\delta\hat{\bw}_3[k]\bigr] \\
\hat{\ba}[k+1]
  &= \hat{\ba}[k]
   + \frac{\ba'_{\mathrm{true}}\bigl(\bw_d[k]\bigr)
           \bigl[\Delta\bw_d[k] + (1-\lambda_c)\,\delta\hat{\bw}_3[k]\bigr]}
          {1 + \dfrac{\ba'_{\mathrm{true}}\bigl(\bw_d[k]\bigr)
           \bigl[\Delta\bw_d[k] + (1-\lambda_c)\,\delta\hat{\bw}_3[k]\bigr]}
          {1-\hat{\ba}[k]}}
\end{align*}

($\hat{\ba}[k]$ now sits in a denominator and $\delta\hat{\bw}_3[k]$ in a
ratio: the line is no longer linear in either)

\textbf{Production and the $b_{\mathrm{true}}$ arm, Euler}

\begin{equation*}
\hat{\ba}[k+1] = \hat{\ba}[k]
  + b\,\bigl(1-\hat{\ba}[k]\bigr)^{2}
    \bigl[\Delta\bw_d[k] + (1-\lambda_c)\,\delta\hat{\bw}_3[k]\bigr]
\end{equation*}

(already nonlinear under Euler: $(1-\hat{\ba}[k])^{2}$ multiplies the
increment. The exact step changes the form of this dependence; it does not
introduce it)

\textbf{The rule that turns a nonlinear line into a mean shift}

\begin{align*}
\E[x\,y] &= \E[x]\,\E[y] + \mathrm{Cov}(x,y) \\
\E\bigl[f(x)\bigr] &= f\bigl(\E[x]\bigr) + \tfrac{1}{2}\,f''\bigl(\E[x]\bigr)\,\mathrm{Var}(x)
  + \mathcal{O}\bigl(\E[(x-\E x)^{3}]\bigr) \\
\E\bigl[f(x,y)\bigr] &= f\bigl(\E x,\E y\bigr)
  + \tfrac{1}{2}\,\frac{\partial^{2}f}{\partial x^{2}}\,\mathrm{Var}(x)
  + \frac{\partial^{2}f}{\partial x\,\partial y}\,\mathrm{Cov}(x,y)
  + \tfrac{1}{2}\,\frac{\partial^{2}f}{\partial y^{2}}\,\mathrm{Var}(y)
  + \cdots
\end{align*}

(for the exact-step oracle line, $x = \hat{\ba}[k]$, $y = \delta\hat{\bw}_3[k]$;
the three second derivatives are those of the Jacobians section of
0902\_formC\_aptrue\_4state.tex, quoted here at
$\Delta\bw_d[k] + (1-\lambda_c)\,\delta\hat{\bw}_3[k] \to 0$)

\begin{align*}
\frac{\partial^{2}\hat{\ba}[k+1]}{\partial\hat{\ba}[k]^{2}}
  &= -\frac{2\,\ba'_{\mathrm{true}}\bigl(\bw_d[k]\bigr)^{2}
        \bigl[\Delta\bw_d[k] + (1-\lambda_c)\,\delta\hat{\bw}_3[k]\bigr]^{2}}
       {\bigl(1-\hat{\ba}[k]\bigr)^{3}} \\
\frac{\partial^{2}\hat{\ba}[k+1]}{\partial\hat{\ba}[k]\,\partial\delta\hat{\bw}_3[k]}
  &= -\frac{2\,(1-\lambda_c)\,\ba'_{\mathrm{true}}\bigl(\bw_d[k]\bigr)^{2}
        \bigl[\Delta\bw_d[k] + (1-\lambda_c)\,\delta\hat{\bw}_3[k]\bigr]}
       {\bigl(1-\hat{\ba}[k]\bigr)^{2}} \\
\frac{\partial^{2}\hat{\ba}[k+1]}{\partial\delta\hat{\bw}_3[k]^{2}}
  &= -\frac{2\,(1-\lambda_c)^{2}\,\ba'_{\mathrm{true}}\bigl(\bw_d[k]\bigr)^{2}}
       {1-\hat{\ba}[k]}
\end{align*}

(all three vanish when the noise is zero, because
$\mathrm{Var}(\hat{\ba})$, $\mathrm{Cov}(\hat{\ba},\delta\hat{\bw}_3)$ and
$\mathrm{Var}(\delta\hat{\bw}_3)$ do; all three scale with the noise power.
This is the whole content of the "cost" of the exact step: the oracle arm,
which under Euler had a linear predict line and no mean shift, acquires these
three terms. Measured on that arm, exact step, seed at truth, 30 seeds:
Meng hold $-0.0037$, canon hold $-0.0022$, amplitude scaling $k^{1.74}$
for $k \in \{0.5, 1, 2\}$; zero noise $-0.00008$ / $-0.00001$.
Which of the three terms carries it, and how the $y_1$ update transforms the
injection into the measured value, is the accounting of Sections
Error dynamics and the probe of 0902\_formC\_aptrue\_4state.tex, not of this note)

\section{An exact step for an exogenous slope that contains no estimate}

(when the slope is a known function of the commanded height, the exact
integral over the step is the difference of the known curve at the two
commanded heights; no $\hat{\ba}$ is needed anywhere in the increment)

\begin{align*}
\hat{\ba}[k+1] &= \hat{\ba}[k]
  + \Bigl[\ba_{\mathrm{true}}\bigl(\bw_d[k+1]\bigr) - \ba_{\mathrm{true}}\bigl(\bw_d[k]\bigr)\Bigr]
  + \ba'_{\mathrm{true}}\bigl(\bw_d[k]\bigr)\,(1-\lambda_c)\,\delta\hat{\bw}_3[k] \\
\ba_{\mathrm{true}}\bigl(\bw_d[k+1]\bigr) - \ba_{\mathrm{true}}\bigl(\bw_d[k]\bigr)
  &= \ba'_{\mathrm{true}}\bigl(\bw_d[k]\bigr)\,\Delta\bw_d[k]
   + \tfrac{1}{2}\,\ba''_{\mathrm{true}}\bigl(\bw_d[k]\bigr)\bigl(\Delta\bw_d[k]\bigr)^{2}
   + \cdots
\end{align*}

(the commanded part of the increment is integrated exactly along the published
curve, so the ratchet of Section 2 is gone; the disturbance part stays first
order in $\delta\hat{\bw}_3$, which is of order $3\times10^{-4}\,R$; the
line is linear in $\hat{\ba}[k]$ and in $\delta\hat{\bw}_3[k]$, so the mean
rule of Section 5 adds nothing. This is the form under which the oracle arm is
both exact and free of a mean shift, and separates the two effects that the
implemented $b_{\mathrm{eff}}$ form mixes: it is not available to production,
whose slope is $b(1-\hat{\ba})^{2}$ by construction)

\end{document}
